\chapter{The model}\label{ch:model}

\section{The impurity}

The impurity is a spin-degenerate orbital with annihilation operators $d_\sigma$,
$\sigma\in\{\up,\dn\}$, occupation $n_d=n_\up+n_\dn$ and spin $\Sz_d=(n_\up-n_\dn)/2$.
Its Hamiltonian is written relative to half filling,
\begin{equation}\label{eq:Hd}
  \Hd=\frac U2\,(n_d-1)^2+\eps\,(n_d-1)+B\,\Sz_d ,
\end{equation}
where $U$ is the charging energy, $\eps=\eps_d+U/2$ is the detuning from the
particle--hole-symmetric point and $B$ is a Zeeman energy that acts on the dot only. In
terms of the conventional Anderson parameters,
$\Hd=\eps_d n_d+Un_\up n_\dn+B\Sz_d+\mathrm{const}$. The four atomic states have the
energies
\begin{equation}\label{eq:atomic}
  E_{\ket0}=\frac U2-\eps,\qquad E_{\ket2}=\frac U2+\eps,\qquad E_{\ket{\sigma}}=\pm\frac B2 ,
\end{equation}
so at $B=0$ the singly occupied doublet is the atomic ground state when
$\abs\eps<U/2$. Multi-orbital impurities with arbitrary one- and two-body terms are
supported as well (\cref{sec:general-impurity}).

\subsection{Charging, gates and the four dot configurations}

The four states in \cref{eq:atomic} are labelled by physical \emph{dot-electron occupations}.
The Andreev-quasiparticle occupations of \cref{sec:abs-occupations} refer to a different basis.
The states $\ket0$ and $\ket2=d_\up\dg d_\dn\dg\ket0$ are spin singlets;
$\ket\up$ and $\ket\dn$ form a doublet. In the conventional Anderson
notation, the first electron-addition energy is $\epsilon_d$ and the second is
$\epsilon_d+U$. A gate that moves these addition energies through the lead
chemical potential changes the preferred charge. The interval
$-U<\epsilon_d<0$, equivalently $|\eps|<U/2$, favours one electron on an
isolated dot. This is the simplest Coulomb-blockade valley.

\begin{figure}[htb]
\centering\manualfigure{primer_dot}
\caption{Physical electron configurations of a single-orbital dot. Tunnelling
visits the singly occupied states by exchanging one electron with a lead.
An induced pair amplitude instead couples the two even configurations directly.
The energy labels use the shifted convention of \cref{eq:Hd}; arrows indicate
couplings between configurations.}
\label{fig:primer-dot}
\end{figure}

Hybridization broadens the charge transitions and mixes dot and lead states.
Near a charge degeneracy, empty or doubly occupied configurations have appreciable
weight: this is a charge-fluctuating or mixed-valence regime. Well inside a
singly occupied valley, with $U$ large relative to the hybridization, charge
fluctuations are mainly virtual and the low-energy variable is a spin. Replacing
the normal leads by superconductors adds coherent pair exchange and removes
their gapless single-particle continuum. These two changes are responsible for
much of the physics that follows.

\begin{physics}[A useful hierarchy of questions]
First ask whether the dot would be empty, singly occupied or doubly occupied in
isolation. Next compare its charge-excitation energies with the coupling to the
leads. Finally ask whether the superconducting gap freezes out lead excitations
before strong spin correlations develop. Identify the regime from the gate
setting together with $U$, $\Gamma$ and $\Delta$.
\end{physics}

\section{Superconducting reservoirs}

Each reservoir $l$ is an $s$-wave BCS superconductor,
\begin{equation}\label{eq:bcs}
  H_l=\sum_{i\sigma}\xi_{li}\,c_{li\sigma}\dg c_{li\sigma}
      -\Delta_l\sum_i\left(c_{li\up}\dg c_{li\dn}\dg+c_{li\dn}c_{li\up}\right),
\end{equation}
with a real gap $\Delta_l>0$; the superconducting phase is moved into the tunnelling terms
(\cref{sec:tunnelling}). The Bogoliubov transformation
\begin{equation}\label{eq:bogoliubov}
  c_{i\up}=u_i\gamma_{i\up}+v_i\gamma_{i\dn}\dg,\qquad
  c_{i\dn}=u_i\gamma_{i\dn}-v_i\gamma_{i\up}\dg,\qquad
  u_i,v_i=\sqrt{\tfrac12\bigl(1\pm\xi_i/E_i\bigr)},
\end{equation}
diagonalizes $H_l$ in terms of quasiparticles with energies $E_i=\sqrt{\xi_i^2+\Delta_l^2}\ge\Delta_l$:
$H_l=\sum_{i\sigma}E_i\gamma_{i\sigma}\dg\gamma_{i\sigma}+\sum_i(\xi_i-E_i)$. The last
term, the energy of the isolated BCS vacuum, is subtracted from all energies reported by
\pkg{}.

\paragraph{The continuum and its discretization.}
A real reservoir has a continuum of levels with a flat normal-state density of states
$\rho=1/(2D)$ per spin on $[-D,D]$; $D$ is the \emph{half}-bandwidth. The impurity
couples to one combination of reservoir states, the contact orbital
\begin{equation}\label{eq:contact}
  C_{l\sigma}=\sum_i\sqrt{w_{li}}\,c_{li\sigma},
\end{equation}
and the impurity couples to the reservoir through the corresponding hybridization function. A finite bath
replaces the continuum by levels $\xi_{li}$ with weights $w_{li}$ that approximate this
function as a quadrature rule. \pkg{} provides logarithmic (\code{cosh_grid}), fitted
surrogate~\cite{baran2023surrogate} and Pad\'e chain-expansion~\cite{bobok2025scalable}
representations. The bath and its weights are part of the model and are stored with every
result. The continuum is recovered only in the limit of many levels; this limit must be
checked separately from all other approximations (\cref{sec:bath,ch:ex-numerics}).

\section{Tunnelling, phases and the reference junction}\label{sec:tunnelling}

The impurity couples to the contact orbitals,
\begin{equation}\label{eq:tunnelling}
  H_T=\sum_{l\sigma}\left(\ee^{\ii\phi_l/2}\,t_l\,C_{l\sigma}\dg d_\sigma+\mathrm{h.c.}\right),
  \qquad \Gamma_l=\pi\rho\,t_l^2,\qquad \Gamma=\sum_l\Gamma_l .
\end{equation}
The phase factor $\ee^{\ii\phi_l/2}$ carries the superconducting phase $\phi_l$ of lead $l$.
In the physical gauge the tunnelling amplitudes are real and the pairing term of lead $l$
is $-\Delta_l\sum_i(\ee^{-\ii\phi_l}c_{li\up}\dg c_{li\dn}\dg+\mathrm{h.c.})$, so that the
pair amplitude $\langle c_{li\up}\dg c_{li\dn}\dg\rangle$ is proportional to
$\ee^{\ii\phi_l}$. The substitution $c_{li\sigma}\to\ee^{-\ii\phi_l/2}c_{li\sigma}$, generated
by the unitary $W=\exp(-\ii\sum_l\phi_lN_l/2)$ with $N_l$ the electron number of lead $l$,
makes the pair potential real, as in \eqref{eq:bcs}, and moves the phase into the
tunnelling. Only phase differences are physical. $\Gamma$ denotes the \emph{total} hybridization. Papers
often quote the coupling $\Gamma_l$ per lead instead, which for two identical leads is
$\Gamma/2$. A direct interlead hopping $W_{LR}$ between the contact orbitals, with the phase
factor $\ee^{\ii(\phi_L-\phi_R)/2}$, models a parallel conduction channel.

The \emph{reference junction} used throughout this manual consists of a single-orbital
impurity, two identical leads with $\Gamma_L=\Gamma_R=\Gamma/2$, phases
$\phi_{L,R}=\mp\phi/2$ and no direct hopping. Its parameters, unless stated otherwise,
are
\begin{equation}\label{eq:reference}
  U=\q{static.u},\quad \Gamma=\q{static.gamma},\quad \eps=\q{static.detuning},\quad
  \phi=\q{static.phi},\quad D=\q{static.bandwidth},\quad \Delta=\q{static.delta},
\end{equation}
with a \code{cosh_grid} bath of $\q{static.pairs}$ level pairs per lead.

\section{Observables and the energy zero}

The supercurrent follows from the phase dependence of the energy,
\begin{equation}\label{eq:current}
  I=\frac{2e}{\hbar}\frac{\partial E}{\partial\phi}=\frac{2e}{\hbar}\left\langle\frac{\partial H}{\partial\phi}\right\rangle ,
\end{equation}
by the Hellmann--Feynman theorem. With $\phi=\phi_R-\phi_L$ and the electron charge $-e$
($e>0$), $I$ is the electric current from the left to the right lead: the tunnelling
changes the electron number of lead $l$ at the rate
$\dot N_l=-(2/\hbar)\langle\partial H/\partial\phi_l\rangle$. The operator $\partial H/\partial\phi$ is available as
the observable \code{phase_derivative}. It is reported in units of $\Delta$, so the
current is in units of $2e\Delta/\hbar$. Further built-in observables are the impurity charge
$n_d$, the impurity spin $\Sz_d$ and the double occupancy $n_\up n_\dn$. Because the BCS
vacuum energies are subtracted, energies at $\Gamma=0$ are the atomic energies
\eqref{eq:atomic}.

The local orbital density matrix resolves occupations and pairing coherences.
In the ordered basis $\ket0,\ket\up,\ket\dn,\ket2$, with
$\ket2=d_\up\dg d_\dn\dg\ket0$, its entries are
$\rho_{ab}=\langle\ket b\!\bra a\rangle$.
Fixed total parity removes local odd--even coherences, while
$\rho_{30}=\langle d_\dn d_\up\rangle=F$ can remain finite.
Writing $P_a=\rho_{aa}$ gives $\langle n_d\rangle=P_\up+P_\dn+2P_2$
and probabilities $(P_0+P_2)/2\pm\operatorname{Re}F$ for the normalized
local states $(\ket0\pm\ket2)/\sqrt2$. Their sum is the local even weight.
Thus unit mean occupation can describe either a local spin or coherent pair mixing.
\code{orbital_density_matrix} evaluates these operators in physical electron
coordinates for QP or DMRG states; \code{orbital_density_properties} also returns
the purity $\operatorname{Tr}\rho^2$. Equal incoherent weighting of a complete
degenerate eigenspace gives a basis-independent density matrix.
\usermanual{docs/interfaces.md} gives the measurement and mixture arguments.

\section{Symmetries and sectors}\label{sec:symmetries}

The Hamiltonian conserves the fermion parity $P=(-1)^N$: pairing changes the electron
number by two. Without transverse fields or spin--orbit coupling it also conserves the
total spin projection $S_z$, and at $B=0$ it is SU(2) symmetric. Calculations are performed in a
\emph{sector} of fixed parity and, optionally, fixed $2S_z$ (\code{Sector(parity, twice_sz)}).
For the symmetric junction at the particle--hole-symmetric point ($\eps=0$, no normal
interlead hopping) there is an additional $\mathbb Z_2$ symmetry of particle--hole type.
It is diagonal in the eta representation (\code{reference_model(symmetry=True)} at $\eps=0$),
and its eigenvalue, the \emph{eta parity} $\eta=\pm1$, can be used as a further
quantum number. Tunnelling, pairing, $U$ and
the phase dependence are eta-even, whereas $n_d-1$ is eta-odd, so a detuning breaks the
symmetry. In \cref{sec:selection} this distinction becomes a selection rule for driven
transitions.

\begin{caution}
At exact degeneracies the eigenvectors can be rotated within their common eigenspace. Expectation values of
non-conserved operators then depend on the basis chosen within the degenerate subspace;
compare complete multiplets, or resolve them by additional quantum numbers.
\end{caution}

\section{Subgap states: singlet and doublet}\label{sec:subgap}

\paragraph{Atomic limit.}
For $\Gamma\to0$ the dot decouples. The spectrum consists of the atomic energies
\eqref{eq:atomic} plus quasiparticle excitations of the leads, each costing at least $\Delta$. Even parity
contains $\ket0$ and $\ket2$, odd parity the doublet $\ket\sigma$.

\paragraph{Large-gap (proximitized-dot) limit.}
In the wide-band limit, when $\Delta$ exceeds the dot and hybridization scales,
virtual Cooper-pair tunnelling induces a local
pairing on the dot. Eliminating the leads gives the effective
Hamiltonian $H_\infty$ of the infinite-gap limit~\cite{meng2009selfconsistent,karrasch2008josephson,zonda2023generalized}
\begin{equation}\label{eq:largegap}
  H_\infty=\Hd-\left(\Gamma_\phi\,d_\up\dg d_\dn\dg+\Gamma_\phi^*\,d_\dn d_\up\right),
  \qquad \Gamma_\phi=\sum_l\Gamma_l\ee^{-\ii\phi_l}=\Gamma\cos\frac\phi2 ,
\end{equation}
the last equality for the symmetric junction. The doublet is unaffected, $E_D=\pm B/2$,
while $\ket0$ and $\ket2$ hybridize into two singlets,
\begin{equation}\label{eq:singlets}
  E_{S,S^*}=\frac U2\mp\sqrt{\eps^2+\abs{\Gamma_\phi}^2}.
\end{equation}
The even subspace is therefore a two-level system with the pseudospin Hamiltonian
$\tfrac U2+\eps\tau_z-\abs{\Gamma_\phi}\tau_x$, where $\tau_z=\pm1$ for $\ket2$ and
$\ket0$. At $B=0$ the ground state is a singlet when $\sqrt{\eps^2+\abs{\Gamma_\phi}^2}>U/2$ and
a doublet otherwise; the phase dependence of $\Gamma_\phi$ drives the transition. This two-level
system reappears in \cref{ch:ex-tls} as the simplest driven model.

\subsection{Diagonalizing the even sector by hand}

In the ordered electron basis $(\ket0,\ket2)$, the even block is
\begin{equation}\label{eq:atomic-even-block}
 H_{\mathrm{even}}=
 \begin{pmatrix}U/2-\eps&-\Gamma_\phi^*\\-\Gamma_\phi&U/2+\eps\end{pmatrix}.
\end{equation}
Its diagonal entries cross when $\eps=0$, but pair mixing turns that crossing
into an avoided crossing of size $2|\Gamma_\phi|$. The lower eigenstate is a
coherent superposition of empty and doubly occupied dot configurations.
For $R=\sqrt{\eps^2+|\Gamma_\phi|^2}>0$, its charge follows either
from its wavefunction or from the Hellmann--Feynman theorem:
\begin{equation}
 \expval{n_d}_S-1=\frac{\partial E_S}{\partial\eps}=-\frac{\eps}{R}.
\end{equation}
At half filling it has mean dot charge one despite containing only empty and
doubly occupied configurations in this limit. This is a direct example of
even parity with odd mean charge.

\begin{figure}[htb]
\centering\manualfigure{primer_atomic}
\caption{Large-gap dot with $U=\q{primer.atomic_u}$ and
$|\Gamma_\phi|=\q{primer.atomic_pairing}$ in a common illustrative energy unit.
(a) The even singlets avoid crossing each other; the lower singlet crosses the
uncoupled odd doublet. (b) Its mean charge and the global ground-state charge.
These curves describe the infinite-gap effective dot. The finite-gap reference
junction is defined separately in \cref{eq:reference}.}
\label{fig:primer-atomic}
\end{figure}

At zero field, a parity transition occurs at $R=U/2$. At fixed
$|\Gamma_\phi|<U/2$, a gate sweep crosses this boundary twice, leaving a
doublet region around half filling. Increasing $|\Gamma_\phi|$ eventually
removes that region. For two leads,
\begin{equation}
 |\Gamma_\phi|^2=\Gamma_L^2+\Gamma_R^2+
                     2\Gamma_L\Gamma_R\cos\phi.
\end{equation}
Pair exchange is strongest when the condensate contributions add, and weakest
when they oppose each other. Exact cancellation at $\phi=\pi$ requires equal
couplings. The even states repel because pairing couples them; a singlet and
doublet cross because their parities differ.

The wide-band qualification is important when using a finite-band solver.
For a flat band of half-width $D$, a lead's zero-energy anomalous contribution
has magnitude
\begin{equation}\label{eq:finite-band-pairing}
 \Gamma_{l,\mathrm{ind}}=
 \frac{2\Gamma_l}{\pi}\arctan\frac{D}{\Delta}.
\end{equation}
This follows by integrating
$\Gamma_l\Delta/[\pi(\xi^2+\Delta^2)]$ over $[-D,D]$.
Taking $D\gg\Delta$ first recovers $\Gamma_l$. Sending $\Delta$ to infinity
at fixed $D$ instead suppresses this contribution. Obtaining the conventional atomic
pairing model therefore requires the wide-band limit first. Replacing each lead
by one paired site defines a separate finite-reservoir approximation.

\subsection{Finite gap and Kondo correlations}
For $\Delta$ comparable to $U$ and $\Gamma$, quasiparticle states of the leads contribute
to the bound states. In the strongly interacting regime, Kondo screening, characterized by
the normal-state Kondo temperature $T_K$, competes with the gap. In the
local-moment regime, $k_BT_K\gg\Delta$ favours a screened singlet and
$k_BT_K\ll\Delta$ a doublet; the boundary depends on conventions, phase and
other parameters~\cite{choi2004kondo,siano2004josephson,bauer2007spectral,meng2009selfconsistent}.
\Cref{ch:correlations} develops this competition and distinguishes it from local pairing.
Besides $S$ and $D$, an excited singlet $S^*$ can lie inside the gap. The continuum starts
one quasiparticle above the ground state, at $E_{\mathrm{ground}}+\Delta$, when
all parity sectors are considered in the continuum limit. In a fixed sector,
the lowest allowed threshold can be higher: a singlet plus one quasiparticle
has odd parity, whereas a doublet plus one quasiparticle can have even parity.
For example, the even breakup threshold on the singlet side can be
$E_D+\Delta$, corresponding to a doublet plus one lead quasiparticle. A finite bath replaces continua by discrete
levels; the shading below indicates the continuum-limit lower envelope.

\begin{figure}[htb]
\centering
\manualfigure{static}
\caption{Subgap spectrum of the reference junction \eqref{eq:reference}, measured from the
ground state: (a) versus the phase at $\eps=\q{static.detuning}$, (b) versus the detuning at
$\phi=\q{static.phi}$. The doublet becomes the ground state near $\phi=\pi$, where
$\abs{\Gamma_\phi}$ is smallest. The shaded region marks the lowest continuum
threshold across parity sectors, $\Delta$ above the ground state. Each probe's
threshold follows from its allowed final sector. QP cutoff $\q{static.cutoff}$,
$\q{static.pairs}$ bath level pairs per lead.}
\label{fig:static}
\end{figure}

\begin{table}[htb]
\centering\small
\caption{Subgap states of the reference junction and their quantum numbers.}\label{tab:states}
\begin{tabularx}{\textwidth}{llllX}
\toprule
State & Parity & Spin & $S_z$ sectors & Character \\
\midrule
$S$ & even & 0 & 0 & bonding combination of $\ket0$, $\ket2$ and lead quasiparticles; ground state at $\phi=\q{static.phi}$ \\
$S^*$ & even & 0 & 0 & antibonding partner of $S$; reached from $S$ by the pair transition \\
$D_{\up,\dn}$ & odd & $\tfrac12$ & $\pm\tfrac12$ & local moment, partially screened; split by $B$ \\
continuum & both & & & ground state plus at least one quasiparticle, energy $\ge E_{\mathrm{ground}}+\Delta$ \\
\bottomrule
\end{tabularx}
\end{table}

\Cref{fig:static} shows the subgap spectrum of the reference junction. At
$\phi=\q{static.phi}$ and $\eps=\q{static.detuning}$ the singlet $S$ is the ground state, the
doublet lies $E_D-E_S=\q{static.doublet_energy}$ above it, and the excited singlet lies at
$E_{S^*}-E_S=\q{static.excited_energy}$, well below the
continuum. The $S$--$S^*$ pair transition and the doublet are the main subjects of the
driven examples.

\begin{physics}
The singlet ground state is a coherent superposition of the empty and doubly occupied dot
(and, at finite $\Delta$, of states with quasiparticles in the leads), held together by
the induced pairing $\Gamma_\phi$. The excited singlet $S^*$ is the antibonding partner of
the same pair. The transition $S\to S^*$ changes the charge distribution while preserving
parity, so a gate or a phase drive can excite it when the matrix element is nonzero.
The doublet carries total spin
$1/2$ and is reached only by changing the parity, that is, by adding or removing a
quasiparticle.
\end{physics}

\section{From a parity crossing to Josephson behaviour}\label{sec:zero-pi-picture}

At zero temperature, the global ground-state energy is the lowest energy across
all relevant sectors. A 0 junction has its equilibrium energy minimum at phase
zero; a $\pi$ junction has its minimum at $\pi$. For an approximately sinusoidal
current, these correspond to positive and negative coefficients of $\sin\phi$.
These definitions classify the full phase dependence of the energy.

In the usual single-level superconducting Anderson model, the singlet branch
tends to favour 0 behaviour and the doublet branch $\pi$ behaviour. Virtual
finite-gap tunnelling processes through a local moment can reverse the pair
transfer amplitude~\cite{glazman1989resonant,spivak1991negative}. The strict
large-gap model instead has a phase-independent doublet energy and therefore
zero doublet current. It demonstrates the parity crossing; reproducing the
finite-gap $\pi$ current requires the virtual processes described above.

\begin{figure}[htb]
\centering\manualfigure{primer_current}
\caption{State-resolved and global ground-state currents of the finite-bath
reference junction, rendered from the same calculation as \cref{fig:static}.
Current is in units of $2e\Delta/\hbar$. The global ground state follows the
singlet or doublet branch with lower energy; disconnected segments indicate
parity changes, where a unique zero-temperature derivative need not exist.
The finite-gap doublet current is distinct from its zero value in the strict
large-gap model. Bath and QP truncations are those of \cref{fig:static}.}
\label{fig:primer-current}
\end{figure}

If the ground-state parity changes as the phase is varied, the ground-state current is
piecewise defined. Both zero and $\pi$ can be local energy minima in intermediate
regimes, conventionally called $0'$ or $\pi'$ according to the global minimum.
Multilevel structure, spin--orbit coupling and magnetic fields admit further
possibilities. Thus ``even means 0'' and ``odd means $\pi$'' are useful
associations for that single-level model. Classify a general weak link by the
phase dependence of its ground-state energy.

At finite temperature and with equilibrium exchange between sectors, use the
free energy $F=-k_BT\ln\sum_m\exp[-E_m/(k_BT)]$ and
$I=(2e/\hbar)\partial_\phi F$. If parity remains trapped during the experiment,
restrict the state sum accordingly. A state-resolved current, a zero-temperature
ground-state current and a thermally averaged current are three different
observables even for the same Hamiltonian.

\begin{physics}[Read the atomic-dot plot]
At half filling, the lower even state in \cref{fig:primer-atomic} has mean charge
one. Which electron configurations produce this mean charge? The state is an
even superposition of zero and two dot electrons. Why can it cross the doublet?
The two states have different parity. What produces a doublet supercurrent at
finite gap? Virtual tunnelling processes generate the phase dependence retained
in \cref{fig:primer-current}.
\end{physics}

\Cref{ch:practical} turns this model into a complete calculation, including
sector comparison, observables, transitions, numerical checks and saved results.

\section{General impurities and models}\label{sec:general-impurity}

\code{make_model} builds models with several reservoirs, arbitrary tunnelling matrices
$T_l$ (including spin mixing), interlead hopping and multi-orbital impurities, specified
either by one- and two-body integrals or directly as fermionic operators. The reference
junction is the special case built by \code{reference_model}. Any Hermitian,
parity-conserving fermionic operator (\code{FermionOperator}) can be turned into a
\code{Hamiltonian} and solved; the examples in \cref{ch:ex-tls,ch:ex-esr} use this freedom.

\tutorial{01_subgap_states}
